\documentclass[12pt]{article}

\usepackage{graphicx}
\usepackage[hypertexnames=false,colorlinks=true,breaklinks]{hyperref}
\usepackage{tikz}
\usepackage{verbatim}
\usepackage{natbib}
\usepackage{apalike}
\usepackage{amsmath,amssymb}

\newcommand{\estNumber}{99331617}

\title{Detecting replay on Emmett's data with svGPFA}
\author{Joaqu\'{i}n Rapela\thanks{j.rapela@ucl.ac.uk}}

\begin{document}

\maketitle

\section{Introduction}

Here I use svGPFA~\citep{dunckerAndSahani18} to explore the low dimensional
structure of Emmett's switching-task recordings.

\section{Methods}

\subsection{Data}

We analyzed the spiking activity of all striatal neurons (neuron indices 0-207)
in electrophysiological recording
\texttt{EJT178\_implant1/recording6\_29\-03\-2022}.
%
Figure~\ref{fig:spikeTimes} plots a sample of spike times used in this
analysis.

We epoched the recordings from 0.2 seconds before to mouse entered the first
port (port 2) to 0.2 seconds after it exited the last port (port 7).
%
Figures~\ref{fig:spikeTimesTrialAt4311.70}-\ref{fig:spikeTimesTrialAt6063.01}
and
figures~\ref{fig:spikeTimesTrialAt6742.00}-\ref{fig:spikeTimesTrialAt7238.00}
plot spike times from behavioral and non-behavioral trials, respectively.


\begin{figure}
    \centering
    \href{http://www.gatsby.ucl.ac.uk/~rapela/emmett/reports/svgpfa_20260307/figures/spike_times_start4045.00_end4345.00.html}{\includegraphics[width=6in]{../../../figures/EJT178_implant1/recording6_29-03-2022/spike_times_start4045.00_end4345.00.png}}

    \caption{Example spikes times used in this analysis. We used recording
    session \texttt{EJT178\_implant1/recording6\_29-03-2022} and neurons
    indices 0 to 207 in the striatum.
    %
    Click on the image to get its interactive version.
    %
    }

    \label{fig:spikeTimes}
\end{figure}

\begin{figure}
    \centering
    \href{http://www.gatsby.ucl.ac.uk/~rapela/emmett/reports/svgpfa_20260307/figures/43952638_epoched_spikes_times_trialID20059.html}{\includegraphics[width=6in]{../../../../svGPFA_striatum/figures/EJT178_implant1/recording6_29-03-2022/43952638_epoched_spikes_times_trialID20059.png}}

    \caption{Example spikes times from the behavioral trial at time 4311.70~sec, from recording
    session \texttt{EJT178\_implant1/recording6\_29-03-2022}.
    %
    Click on the image to get its interactive version.
    %
    }

    \label{fig:spikeTimesTrialAt4311.70}
\end{figure}

\begin{figure}
    \centering
    \href{http://www.gatsby.ucl.ac.uk/~rapela/emmett/reports/svgpfa_20260307/figures/43952638_epoched_spikes_times_trialID20061.html}{\includegraphics[width=6in]{../../../../svGPFA_striatum/figures/EJT178_implant1/recording6_29-03-2022/43952638_epoched_spikes_times_trialID20061.png}}

    \caption{Example spikes times from the behavioral trial at time 4320.08~sec, from recording
    session \texttt{EJT178\_implant1/recording6\_29-03-2022}.
    %
    Click on the image to get its interactive version.
    %
    }

    \label{fig:spikeTimesTrialAt4320.08}
\end{figure}

\begin{figure}
    \centering
    \href{http://www.gatsby.ucl.ac.uk/~rapela/emmett/reports/svgpfa_20260307/figures/43952638_epoched_spikes_times_trialID20084.html}{\includegraphics[width=6in]{../../../../svGPFA_striatum/figures/EJT178_implant1/recording6_29-03-2022/43952638_epoched_spikes_times_trialID20084.png}}

    \caption{Example spikes times from the behavioral trial at time 4448.36~sec, from recording
    session \texttt{EJT178\_implant1/recording6\_29-03-2022}.
    %
    Click on the image to get its interactive version.
    %
    }

    \label{fig:spikeTimesTrialAt4448.36}
\end{figure}

\begin{figure}
    \centering
    \href{http://www.gatsby.ucl.ac.uk/~rapela/emmett/reports/svgpfa_20260307/figures/43952638_epoched_43952638_epoched_spikes_times_trialID20378.html}{\includegraphics[width=6in]{../../../../svGPFA_striatum/figures/EJT178_implant1/recording6_29-03-2022/43952638_epoched_spikes_times_trialID20378.png}}

    \caption{Example spikes times from the behavioral trial at time 6063.01~sec, from recording
    session \texttt{EJT178\_implant1/recording6\_29-03-2022}.
    %
    Click on the image to get its interactive version.
    %
    }

    \label{fig:spikeTimesTrialAt6063.01}
\end{figure}

\begin{figure}
    \centering
    \href{http://www.gatsby.ucl.ac.uk/~rapela/emmett/reports/svgpfa_20260307/figures/23323766_pseudo_epoched_spikes_times_trialID30185.html}{\includegraphics[width=6in]{../../../../svGPFA_striatum/figures/EJT178_implant1/recording6_29-03-2022/23323766_pseudo_epoched_spikes_times_trialID30185.png}}

    \caption{Example spikes times from the non-behavioral trial at time 6742.00~sec, from recording
    session \texttt{EJT178\_implant1/recording6\_29-03-2022}.
    %
    Click on the image to get its interactive version.
    %
    }

    \label{fig:spikeTimesTrialAt6742.00}
\end{figure}

\begin{figure}
    \centering
    \href{http://www.gatsby.ucl.ac.uk/~rapela/emmett/reports/svgpfa_20260307/figures/23323766_pseudo_epoched_spikes_times_trialID30309.html}{\includegraphics[width=6in]{../../../../svGPFA_striatum/figures/EJT178_implant1/recording6_29-03-2022/23323766_pseudo_epoched_spikes_times_trialID30309.png}}

    \caption{Example spikes times from the non-behavioral trial at time 7238.00~sec, from recording
    session \texttt{EJT178\_implant1/recording6\_29-03-2022}.
    %
    Click on the image to get its interactive version.
    %
    }

    \label{fig:spikeTimesTrialAt7238.00}
\end{figure}

\subsection{svGPFA model}

The svGPFA model estimates $K$ latent variables in the form of Gaussian
processes:

\begin{align*}
    x_k^{(r)}(t)\sim GP(\mu_k(t), \kappa_k(t,t'))\quad k=1,\ldots,K
\end{align*}

\noindent I concatenate the latent variables $x_k^{(r)}(t)$ into a latent vector
$\mathbf{x}^{(r)}(t)$:

\begin{align}
    \mathbf{x}^{(r)}(t)=\left[\begin{array}{c}
        x_1^{(r)}(t)\\
                            \vdots\\
                            x_K^{(r)}(t)
                        \end{array}\right]
    \label{eq:latentsVector}
\end{align}

The latent variables are combined linearly to produced the log pre-intensity,
$\mathbf{h}^{(r)}(t)=[h_1^{(r)}(t),\ldots,h_N^{(r)}(t)]^\intercal$, where
$h_n^{(r)}(t)$ is the log
pre-intensity of neuron $n$ and trial $r$, for a total of $N$ neurons:

\begin{align}
    \mathbf{h}^{(r)}(t)=C\mathbf{x}^{(r)}(t)+\mathbf{d}
    \label{eq:logPreIntensity}
\end{align}

The log pre-intensity of neuron $n$ and trial $r$ is exponentiated to produce the conditional
intensity function of neuron $n$ and trial $r$:

\begin{align}
    \lambda_n^{(r)}(t)=\exp(h_n^{(r)}(t))
    \label{eq:cif}
\end{align}

\noindent which can be interpreted as the instantaneous firing rate of neuron
$n$ and trial $r$ at time $t$.

The conditional intensity function of neuron $n$ in trial $r$ is used to calculate the
probability of spikes $\{t_{i,n}\}_{i=1}^{\Phi(n,r)}$ in a trial $r$ of length $\tau$:

\begin{align*}
    P(\{t_{i,n}\}_{i=1}^{\Phi(n,r)}|\lambda_n^{(r)}(t))=\exp\left(-\int_0^\tau\lambda_n^{(r)}(t)dt\right)\prod_{i=1}^{\Phi(n,r)}\lambda_n^{(r)}(t_{i,n})
\end{align*}

\section{Results}

\listoffigures

\begin{figure}
    \centering
    \href{http://www.gatsby.ucl.ac.uk/~rapela/emmett/reports/svgpfa_20260307/figures/54368807_orthonormalized_nonEpoched_latents.html}{\includegraphics[width=6in]{../../../figures/EJT178_implant1/recording6_29-03-2022/54368807_orthonormalized_nonEpoched_latents.png}}
    \caption{Orthonormalised latent variables ($\mathbf{x}(t)$ in
    Eq.~\ref{eq:latentsVector}) inferred for data in the behavioral period
    trials used to estimate the svGPFA model parameters (i.e., trials 42-141).
    %
    Trials have been concatenated in continuous time.
    %
    The colored vertical lines indicate `poke\_in` times, with
    green, red, blue, yellow, purple, cyan and magenta corresponding to ports
    1 to 7, respectively.
    %
    The correct sequence is $2\rightarrow 1\rightarrow 6\rightarrow 3\rightarrow 7$
    corresponding to
    $\text{red}\rightarrow\text{green}\rightarrow\text{blue}\rightarrow\text{cyan}\rightarrow\text{magenta}$.
    %
    Click on the image to get its interactive version.
    %
    }
    \label{fig:inferredLatentsTrain}
\end{figure}

\begin{figure}
    \centering

    \href{http://www.gatsby.ucl.ac.uk/~rapela/emmett/reports/svgpfa_20260307/figures/54368807_hmm_mean.html}{\includegraphics[width=6in]{../../../figures/EJT178_implant1/recording6_29-03-2022/54368807_hmm_mean.png}}

    \caption{HMM emission parameters. Inferred means and standard deviations (error
    bars) of the Gaussian emission model for each latent (i.e., observation)
    and for each HMM state.}

    \label{fig:hmm_emision_params}
\end{figure}

\begin{figure}
    \centering

    \href{http://www.gatsby.ucl.ac.uk/~rapela/emmett/reports/svgpfa_20260307/figures/54368807_hmm_transition.html}{\includegraphics[width=6in]{../../../figures/EJT178_implant1/recording6_29-03-2022/54368807_hmm_transition.png}}

    \caption{HMM transition matrix.}

    \label{fig:hmm_transition_matrix}
\end{figure}

\begin{figure}
    \centering

    \href{http://www.gatsby.ucl.ac.uk/~rapela/emmett/reports/svgpfa_20260307/figures/train54368807_test54368807_true_vs_inferred_hmm_most_prob_state_seq.html}{\includegraphics[width=6in]{../../../figures/EJT178_implant1/recording6_29-03-2022/train54368807_test54368807_true_vs_inferred_hmm_most_prob_state_seq.png}}

    \caption{True (blue) versus estimated (red) behavioral states with
    estimation and inference performed on the same data segment of the
    behavioral period. The svGPFA and
    HMM models were estimated with data from trials 42-141. The svGPFA model
    inferred latents and the HMM inferred behavioral states from the same
    trials.}

    \label{fig:trueVsInferredStatesTest42-141}
\end{figure}

\begin{figure}
    \centering
    \href{http://www.gatsby.ucl.ac.uk/~rapela/emmett/reports/svgpfa_20260307/figures/92418550_orthonormalized_nonEpoched_latents.html}{\includegraphics[width=6in]{../../../figures/EJT178_implant1/recording6_29-03-2022/92418550_orthonormalized_nonEpoched_latents.png}}

    \caption{Orthonormalised latent variables inferred from neural data in the
    behavioral period not used to train svGPFA model parameters (i.e., trials
    242-341).}

    \label{fig:inferredLatentsTestDuringBehavior}
\end{figure}

\begin{figure}
    \centering

    \href{http://www.gatsby.ucl.ac.uk/~rapela/emmett/reports/svgpfa_20260307/figures/train54368807_test92418550_true_vs_inferred_hmm_most_prob_state_seq.html}{\includegraphics[width=6in]{../../../figures/EJT178_implant1/recording6_29-03-2022/train54368807_test92418550_true_vs_inferred_hmm_most_prob_state_seq.png}}

    \caption{True (blue) versus estimated (red) behavioral states, when
    estimation and inference were done on different data segments of the
    behavioral period. The svGPFA and HMM models were estimated with data in
    trials 42-141. The svGPFA model inferred latents and the HMM inferred
    behavioral states from data in trials 242-341.}

    \label{fig:trueVsInferredStatesTest242-341}
\end{figure}

\begin{figure}
    \centering

    \href{http://www.gatsby.ucl.ac.uk/~rapela/emmett/reports/svgpfa_20260307/figures/87796368_orthonormalized_nonEpoched_latents.html}{\includegraphics[width=6in]{../../../figures/EJT178_implant1/recording6_29-03-2022/87796368_orthonormalized_nonEpoched_latents.png}}
    \caption{Orthonormalised latent variables inferred from data in a
    non-behavioral between 6,000 and 6,400 seconds.}

    \label{fig:inferredLatentsTestDuringPostBehavior6000_6400}
\end{figure}

\begin{figure}
    \centering

    \href{http://www.gatsby.ucl.ac.uk/~rapela/emmett/reports/svgpfa_20260307/figures/train54368807_test87796368_inferred_hmm_most_prob_state_seq.html}{\includegraphics[width=6in]{../../../figures/EJT178_implant1/recording6_29-03-2022/train54368807_test87796368_inferred_hmm_most_prob_state_seq.png}}

    \caption{True (blue) versus estimated (red) behavioral states, when
    estimation was done with data in the behavioral period and inference was
    done with data in a non-behavioral period between 6,000 and 6,400 seconds.}

    \label{fig:trueVsInferredStatesTest6000_6400}
\end{figure}

\begin{figure}
    \centering

    \href{http://www.gatsby.ucl.ac.uk/~rapela/emmett/reports/svgpfa_20260307/figures/71005668_orthonormalized_nonEpoched_latents.html}{\includegraphics[width=6in]{../../../figures/EJT178_implant1/recording6_29-03-2022/71005668_orthonormalized_nonEpoched_latents.png}}
    \caption{Orthonormalised latent variables inferred from data in a
    non-behavioral between 6,400 and 6,800 seconds.}

    \label{fig:inferredLatentsTestDuringPostBehavior6400_6800}
\end{figure}

\begin{figure}
    \centering

    \href{http://www.gatsby.ucl.ac.uk/~rapela/emmett/reports/svgpfa_20260307/figures/train54368807_test71005668_inferred_hmm_most_prob_state_seq.html}{\includegraphics[width=6in]{../../../figures/EJT178_implant1/recording6_29-03-2022/train54368807_test71005668_inferred_hmm_most_prob_state_seq.png}}

    \caption{True (blue) versus estimated (red) behavioral states, when
    estimation was done with data in the behavioral period and inference was
    done with data in a non-behavioral period between 6,400 and 6,800 seconds.}

    \label{fig:trueVsInferredStatesTest6400_6800}
\end{figure}

\begin{figure}
    \centering

    \href{http://www.gatsby.ucl.ac.uk/~rapela/emmett/reports/svgpfa_20260307/figures/07996538_orthonormalized_nonEpoched_latents.html}{\includegraphics[width=6in]{../../../figures/EJT178_implant1/recording6_29-03-2022/07996538_orthonormalized_nonEpoched_latents.png}}
    \caption{Orthonormalised latent variables inferred from data in a
    non-behavioral between 6,800 and 7,200 seconds.}

    \label{fig:inferredLatentsTestDuringPostBehavior6800_7200}
\end{figure}

\begin{figure}
    \centering

    \href{http://www.gatsby.ucl.ac.uk/~rapela/emmett/reports/svgpfa_20260307/figures/train54368807_test07996538_inferred_hmm_most_prob_state_seq.html}{\includegraphics[width=6in]{../../../figures/EJT178_implant1/recording6_29-03-2022/train54368807_test07996538_inferred_hmm_most_prob_state_seq.png}}

    \caption{True (blue) versus estimated (red) behavioral states, when
    estimation was done with data in the behavioral period and inference was
    done with data in a non-behavioral period between 6,800 and 7,200 seconds.}

    \label{fig:trueVsInferredStatesTest6800_7200}
\end{figure}

\begin{figure}
    \centering

    \href{http://www.gatsby.ucl.ac.uk/~rapela/emmett/reports/svgpfa_20260307/figures/96281561_orthonormalized_nonEpoched_latents.html}{\includegraphics[width=6in]{../../../figures/EJT178_implant1/recording6_29-03-2022/96281561_orthonormalized_nonEpoched_latents.png}}
    \caption{Orthonormalised latent variables inferred from data in a
    non-behavioral between 7,200 and 7,600 seconds.}

    \label{fig:inferredLatentsTestDuringPostBehavior7200_7600}
\end{figure}

\begin{figure}
    \centering

    \href{http://www.gatsby.ucl.ac.uk/~rapela/emmett/reports/svgpfa_20260307/figures/train54368807_test96281561_inferred_hmm_most_prob_state_seq.html}{\includegraphics[width=6in]{../../../figures/EJT178_implant1/recording6_29-03-2022/train54368807_test96281561_inferred_hmm_most_prob_state_seq.png}}

    \caption{True (blue) versus estimated (red) behavioral states, when
    estimation was done with data in the behavioral period and inference was
    done with data in a non-behavioral period between 7,200 and 7,600 seconds.}

    \label{fig:trueVsInferredStatesTest7200_7600}
\end{figure}

\bibliographystyle{apalike}
\bibliography{latentVariablesModels}

\end{document}
